\section*{Chapter \ref{ch:fm:commercial-relationships-with-venues-brokers-and-vendors} --- Commercial Relationships with Venues, Brokers and Vendors}

\begin{solution}{exo:fm:commercial-relationships-with-venues-brokers-and-vendors:1}
$150$ million shares $\times\$0.000075=\$11\,250$ a day, \$2.83 million over 252 days.
\end{solution}

\begin{solution}{exo:fm:commercial-relationships-with-venues-brokers-and-vendors:2}
The best alternative is another clearing broker (or self-clearing); the \omterm{def:fm:commercial-relationships-with-venues-brokers-and-vendors:batna}{reservation value} is the all-in cost of that alternative, including
margin terms and the cost of moving: the worst terms the firm should accept from its current broker.
\end{solution}

\begin{solution}{exo:fm:commercial-relationships-with-venues-brokers-and-vendors:3}
It pads 102 million shares a day at \$0.001, \$102\,000, and quotes 1\,620 extra symbols at \$2, \$3\,240.
\end{solution}

\begin{solution}{exo:fm:commercial-relationships-with-venues-brokers-and-vendors:4}
See the proof of \cref{prop:fm:commercial-relationships-with-venues-brokers-and-vendors:nash}. Large firm: $C=0$, $r_f=10\,000$, $G=36\,000$,
$\beta=0.5$: $T^*=10\,000+0.5\times26\,000=\$23\,000$ a day, \$0.000096 a share on 240 million shares.
\end{solution}

\begin{solution}{exo:fm:commercial-relationships-with-venues-brokers-and-vendors:5}
Its costs of meeting the obligations, \$105\,240 a day, exceed the venue's gain from its liquidity, \$22\,500: no transfer satisfies both. Tailored
terms remove the costs, and the venue's \$7\,200 gain is shared.
\end{solution}

\begin{solution}{exo:fm:commercial-relationships-with-venues-brokers-and-vendors:6}
Its volume is most uncertain then, so the probability and size of a shortfall are largest; the commitment is a put it has sold on its own growth.
\end{solution}

\begin{solution}{exo:fm:commercial-relationships-with-venues-brokers-and-vendors:7}
$10\,000+0.3\times26\,000=\$17\,800$ a day, \$0.000074 a share: slightly below the published rate.
\end{solution}

\begin{solution}{exo:fm:commercial-relationships-with-venues-brokers-and-vendors:8}
\$2.8 million a year is the rebate at the required volume; a firm below it must pad its volume and quote extra symbols to qualify, which on the
chapter's costs costs more than the rebate unless it is already above 1.16\% of consolidated volume.
\end{solution}

\begin{solution}{pb:fm:commercial-relationships-with-venues-brokers-and-vendors:1}
\begin{enumerate}
\item See \cref{def:fm:commercial-relationships-with-venues-brokers-and-vendors:batna,def:fm:commercial-relationships-with-venues-brokers-and-vendors:zopa}.
\item See \cref{prop:fm:commercial-relationships-with-venues-brokers-and-vendors:nash}.
\item Venue: programmes, obligations' measurement, connectivity; broker: rates, margin and financing, service; vendor: licences, counts, terms.
\item See \cref{def:fm:commercial-relationships-with-venues-brokers-and-vendors:commit}.
\item See \cref{dat:fm:commercial-relationships-with-venues-brokers-and-vendors:qmm}.
\item Padded shares at a loss per share and extra symbol-days at the NBBO at a cost per symbol.
\item Negative below the requirement, positive from 1.16\% of consolidated volume.
\item Padding costs \$0.001 a share against a rebate of \$0.000075, so each padded share costs about thirteen times its rebate.
\item Net $-\$93\,990$ a day; venue gain \$22\,500; empty zone.
\item Zone \$0 to \$7\,200 a day; Nash \$3\,600, \$0.000075 a share.
\item Zone \$10\,000 to \$36\,000; Nash \$23\,000 (\$0.000096 a share) with the outside option, \$18\,000 without.
\item \$5\,000 a day: half of its value, by \cref{prop:fm:commercial-relationships-with-venues-brokers-and-vendors:nash}.
\item At 0.3, \$17\,800 a day (\cref{exo:fm:commercial-relationships-with-venues-brokers-and-vendors:7}).
\item Not on the published terms; it should ask for requirements measured on its own activity.
\item Yes; it should ask for more than the published rebate, using its rival venue as its outside option.
\item Connect to a second venue with a programme sized for small firms.
\item Measurement of obligations, capped penalties, notice and termination rights, most-favoured-nation terms.
\item Monthly against the firm's own activity records, with alerts when a requirement is at risk.
\item Above 1.16\% of consolidated volume the published programme pays; Nash gives the small firm \$0.000075 a share on tailored terms (nothing on
the published ones) and the large firm \$0.000096.
\item Small: decline the published terms and ask for a programme sized to its activity, while building a second venue as an alternative. Large:
join, and ask for a higher rebate backed by its rival venue's offer.
\end{enumerate}
\end{solution}

\begin{solution}{iq:fm:commercial-relationships-with-venues-brokers-and-vendors:1}
A rebate or discount for quoting and volume obligations; the venue gets displayed liquidity that attracts takers, whose fees it earns.
\iqlookfor{liquidity attracts flow.}
\end{solution}

\begin{solution}{iq:fm:commercial-relationships-with-venues-brokers-and-vendors:2}
Another clearing broker, or self-clearing if the firm is large enough; its value sets the \omterm{def:fm:commercial-relationships-with-venues-brokers-and-vendors:batna}{reservation value}.
\iqlookfor{a concrete alternative.}
\end{solution}

\begin{solution}{iq:fm:commercial-relationships-with-venues-brokers-and-vendors:3}
As an option: expected rebate in qualifying months against the cost of padding in months that would otherwise miss, and the value of the cliff on
all shares in qualifying months.
\iqlookfor{expected value over the volume distribution.}
\end{solution}

\begin{solution}{iq:fm:commercial-relationships-with-venues-brokers-and-vendors:4}
Maximise $(T-C-r_f)^\beta(G-T-r_v)^{1-\beta}$; $T^*=C+r_f+\beta(G-C-r_f-r_v)$.
\iqlookfor{outside options shift the split.}
\end{solution}

\begin{solution}{iq:fm:commercial-relationships-with-venues-brokers-and-vendors:5}
Paying penalties when volume falls, being pushed to trade unprofitably to meet it, and lock-in to the counterparty.
\iqlookfor{the sold put.}
\end{solution}

\begin{solution}{iq:fm:commercial-relationships-with-venues-brokers-and-vendors:6}
Compare the all-in value per share at the flow each venue would get, including tier cliffs, obligations and fill quality; the best split may meet
one venue's threshold rather than dividing evenly.
\iqlookfor{cliffs and all-in value.}
\end{solution}
